\clearpage\fancyhead[L]{\small GAUSS AN VON LINDENAU.}\SourceLeaf{321}

\begin{center}\textbf{GAUSS AN VON LINDENAU.}\par\vspace{4pt}\hrule\vspace{4pt}\fontsize{10}{12}\selectfont Monatliche Correspondenz. Bd. XXI. S. 276\,..\,280. März 1810.\par\vspace{4pt}\hrule\vspace{4pt}\end{center}
\begin{flushright}Göttingen, 23. Februar 1810.\end{flushright}

\par\indent ---\enspace---\enspace--- Ich habe gerade einige Rechnungen zu den beiden letzten Oppositionen der Pallas abgeschlossen. Die Ergebnisse habe ich in Nr. 32 der Göttinger gelehrten Zeitungen [S. 317 dieses Bandes] abdrucken lassen und teile sie Ihnen hier mit. ---\enspace---

Ich habe daraus bereits einige weitere Ergebnisse zu den Elementen gewonnen und arbeite jetzt noch an einigen Rechnungen dazu. All dies soll in einen Aufsatz für unsere Sozietät eingehen. Darin werde ich auch verschiedene Kunstgriffe erläutern, die ich seit vielen Jahren bei der Anwendung der im 3. Abschnitt des 2. Buches meiner \emph{Theoria} entwickelten Methode (\emph{des moindres carrés}) verwende. Sobald die Arbeit fertig ist, schicke ich Ihnen einen Auszug.

Ich füge die hiesigen Beobachtungen der \emph{Vesta} bei, die alle Professor \textsc{Harding} am Mauerquadranten gemacht hat.

\Needspace{12\baselineskip}\begin{center}\emph{Beobachtungen der Vesta.}\end{center}
\begin{center}\fontsize{10}{12}\selectfont\setlength{\tabcolsep}{5pt}\renewcommand{\arraystretch}{1.15}\begin{tabular}{lrrr}
1810 & \shortstack{Mittlere Zeit\\in Göttingen} & \shortstack{Scheinbare Rektaszension\\der Vesta} & \shortstack{Scheinbare nördliche\\Deklination}\\
Januar 13 & $11^{\text{u}}\;1^{\text{m}}\;27^{\text{s}}7$ & $98^\circ\;5'\;44''5$ & $23^\circ\;19'\;17''6$\\
15 & $10\;51\;30.1$ & $97\;34\;57.0$ & $23\;26\;14.3$\\
28 & $9\;49\;13.7$ & $94\;46\;15.2$ & $24\;7\;33.7$\\
Februar 2 & $9\;26\;32.7$ & $94\;0\;51.3$ & $24\;20\;51.7$\\
3 & $9\;22\;7.4$ & $93\;53\;19.2$ & $24\;23\;37.8$\\
6 & $9\;8\;56.1$ & $93\;32\;33.2$ & $24\;31\;1.9$\\
\end{tabular}\end{center}

\Needspace{21\baselineskip}\par\indent Hier sind auch unsere Beobachtungen der Bedeckung des Jupiter vom 8. Februar. Die Austritte sind ganz unsicher, weil der Mond sehr tief stand.

\Needspace{17\baselineskip}\begin{center}\emph{Eintritte.}\end{center}
\begin{center}\fontsize{10}{12}\selectfont\setlength{\tabcolsep}{5pt}\renewcommand{\arraystretch}{1.0}\begin{tabular}{lrl}
1. Trabant & $8^{\text{u}}\;58^{\text{m}}\;2^{\text{s}}9$ & \textsc{Harding}\\
 & $3.3$ & \textsc{Gauss}\\
2. Trabant & $8\;58\;26.2$ & G.\\
 & $28.7$ & H.\\
$\JupiterSymbol$ I. R. & $9\;0\;16.4$ & H.\\
 & $17.4$ & G.\\
$\JupiterSymbol$ II. R. & $9\;2\;26.0$ & G.\\
 & $26.7$ & H.\\
3. Trabant & $9\;7\;55.5$ & G.\\
 & $56.4$ & H.\\
4. Trabant & $9\;14\;35.3$ & H.\\
 & $37.0$ & G.\\
\end{tabular}\end{center}
\Needspace{9\baselineskip}\begin{center}\emph{Austritte.}\end{center}
\begin{center}\fontsize{10}{12}\selectfont\setlength{\tabcolsep}{5pt}\renewcommand{\arraystretch}{1.0}\begin{tabular}{lrl}
1. Trabant & $9^{\text{u}}\;30^{\text{m}}\;59^{\text{s}}$ & \textsc{Gauss} schon ausgetreten\\
$\JupiterSymbol$ I. R. & $9\;33\;37.9$ & H.\\
II. R. & $9\;35\;19.6$ & G.\\
 & $36.1$ & H.\\
3. Trabant & $9\;43\;39.3$ & H.\\
\end{tabular}\end{center}

\par\indent Für den Hinweis auf die Druckfehler in meiner Theoria bin ich Herrn \textsc{Oriani} sehr verbunden. Er hat ganz recht, dass ich auf S. 129 vergessen habe hinzuzufügen, dass $B, B', B'' = 0$ vorausgesetzt
\SourceLeaf{322}werden muss, wenn die Bedingungsgleichung, bei der die Gleichung [7] unbrauchbar ist, die dort angegebene Form haben soll. Es ist übrigens klar, dass die Gleichung [7] auch dann unbrauchbar sein kann, wenn $B, B', B''=0$ nicht gilt, nämlich wenn der Ausdruck mit zwölf Gliedern, den \textsc{Oriani} entwickelt hat, zufällig $=0$ oder sehr klein wird. Mir war schon früher eingefallen, dass \textsc{Euler} bereits das Theorem gefunden hatte, aus dem sich der schöne, von mir \textsc{La Place} zugeschriebene Lehrsatz sehr leicht ableiten lässt. Die Stelle auf S. 212 war aber schon gedruckt. Ich wollte es jedoch nicht in die Errata aufnehmen, weil \textsc{La Place} das genannte Theorem wenigstens als Erster in der dort verwendeten Form aufgestellt hat. Die meisten der von \textsc{Oriani} gemeldeten Druckfehler hatte ich mir auch schon notiert. Hier sind noch drei weitere, die er übersehen hat:

\begin{center}\fontsize{10}{12}\selectfont\setlength{\tabcolsep}{5pt}\renewcommand{\arraystretch}{1.15}\begin{tabular}{lrlrlllll}
S. & 1 & Zeile & 4 & v. u. & statt & inversa & lies & composita\\
--- & 65 & --- & 2 & --- & --- & $\cos\theta$. & lies & $\cos\zeta$\\
--- & 195 & --- & 16 &  & --- & $\lambda'''$ & lies & $=\lambda'''$.\\
\end{tabular}\end{center}

\par\indent Den Druckfehler $380^\circ$ statt $180^\circ$ auf S. 4 finde ich in meinem Exemplar nicht. Anstelle der zwei Druckfehler auf S. 83 könnte man einen einzigen angeben, nämlich $\zeta+45$ statt $\zeta-45^\circ$.
